% Volume IV: Law as an Epistemic Machine
% Part VII. Evidence Before Action
% Chapter 35: The Architecture of Proof
% Spine: proof-spines/ch035-spine.md

\chapter{The Architecture of Proof}
\label{ch:ch035}

Legal proof separates factual belief from legally warranted judgment. A tribunal may suspect, even strongly, without yet occupying the institutional position from which action is authorized. The chapter therefore models proof as an \textbf{authorized evidentiary path} rather than as naked confidence:
\[
E
\xrightarrow{\operatorname{admit}}
E^\ast
\xrightarrow{\operatorname{test}}
\widetilde E
\xrightarrow{\operatorname{weigh}}
J
\xrightarrow{\operatorname{authorize}}
A.
\]
\Definition\ Here raw material \(E\) is first admitted into the record as \(E^\ast\), then adversarially tested into \(\widetilde E\), then weighed into judgment \(J\), and only then converted into authorized action \(A\).
Each arrow is a distinct site of institutional discipline and a possible refusal point. Legal proof is not belief with more confidence added; it is confidence produced through this public chain.
